\documentclass[11pt]{article}
\usepackage[margin=1in]{geometry}
\usepackage{amsmath,amssymb}
\usepackage{booktabs}
\usepackage{listings}
\usepackage[hidelinks]{hyperref}

\lstset{basicstyle=\ttfamily\footnotesize,breaklines=true,
        columns=fullflexible,frame=single,framesep=4pt,xleftmargin=6pt}

\newcommand{\WL}{Wolfram Language}

\title{Closed-form solution of the $\xi$ separated equation:\\
       what three computer algebra systems can and cannot do}
\author{Brent Baccala}
\date{23 August 2026}

\begin{document}
\maketitle

\section{The equation}

The separation of the hydrogenic problem in parabolic coordinates
$\xi=r+z$, $\eta=r-z$, $\psi=\tan^{-1}(y/x)$ produces, for the $\xi$
factor,
\begin{equation}\label{eq:f1}
\frac{\mathrm{d}}{\mathrm{d}\xi}\left(\xi\frac{\mathrm{d}f_1}{\mathrm{d}\xi}\right)
  +\left[\tfrac{1}{2}E\xi-\tfrac{1}{4}m^2/\xi+\beta_1\right]f_1 = 0 .
\end{equation}
Everything other than $\xi$ and $f_1$ is constant: the energy $E$, the
magnetic quantum number $m$, and the separation constant $\beta_1$.
Note that $\xi = r+z \geq 0$ identically, since $r\geq|z|$.

Multiplying through by $\xi$ puts \eqref{eq:f1} in the form in which its
character is visible,
\begin{equation}\label{eq:f1x}
\xi^2 f_1'' + \xi f_1' + \left(\tfrac{1}{2}E\xi^2 + \beta_1\xi
  - \tfrac{1}{4}m^2\right) f_1 = 0 ,
\end{equation}
which is \emph{not} Bessel's equation unless $\beta_1=0$: the term
$\beta_1\xi$ is linear in $\xi$ where Bessel's equation requires a pure
$k^2\xi^2$. Substituting $f_1 = \xi^{-1/2}\,w(2k\xi)$ with
$k=\sqrt{-E/2}$ converts \eqref{eq:f1x} into Whittaker's equation
\begin{equation}\label{eq:whittaker}
\frac{\mathrm{d}^2 w}{\mathrm{d}z^2}
  + \left(-\frac{1}{4} + \frac{\kappa}{z}
          + \frac{\tfrac{1}{4}-\mu^2}{z^2}\right) w = 0 ,
\qquad \kappa = \frac{\beta_1}{2k},\quad \mu = \frac{m}{2},
\end{equation}
so the general solution is
$f_1 = \xi^{-1/2} W_{\kappa,\,m/2}(2k\xi)$ and its $M$ companion.  The
separation constant $\beta_1$ is the Coulomb (Sommerfeld) parameter
$\kappa$; setting $\beta_1=0$ sends $\kappa\to0$ and only then does
\eqref{eq:f1x} degenerate to Bessel.

\section{Summary of results}

Four systems were tried on \eqref{eq:f1} with all of $E$, $m$, $\beta_1$
symbolic.  Versions: Wolfram Engine 14.0.0 (21 December 2023), SymPy
1.14.0, SageMath 10.7 driving Maxima 5.48.1, and FriCAS 1.3.13 on SBCL
2.4.8.  Numerical checks throughout use mpmath 1.3.0.

\begin{table}[h]
\centering
\small
\setlength{\tabcolsep}{4pt}
\begin{tabular}{@{}llll@{}}
\toprule
 & \textbf{general $\beta_1\neq0$} & \textbf{$\beta_1=0$} & \textbf{closed form returned}\\
\midrule
Wolfram 14.0        & solved & solved & \verb|HypergeometricU|, \verb|LaguerreL| / \verb|BesselJ|,\verb|BesselY| \\
Sage 10.7 / Maxima  & solved\rlap{$^\dagger$} & solved & \verb|hypergeometric_M|,\verb|_U| / \verb|bessel_J|,\verb|bessel_Y| \\
SymPy 1.14.0        & \textbf{no} & solved & --- / \verb|besselj|, \verb|bessely| \\
FriCAS 1.3.13       & \textbf{no}\rlap{$^\ddagger$} & odd $m$ only\rlap{$^\ddagger$} & Liouvillian basis, else \verb|[]| \\
\bottomrule
\end{tabular}
\end{table}

\noindent$^\dagger$ only with \verb|contrib_ode=True|, and on a branch that
degenerates at integer $m$; see \S\ref{sec:sage}.
$^\ddagger$ FriCAS is not searching a special-function table at all; it
\emph{decides} Liouvillian solvability, and answers a sharper question than
the other three.  See \S\ref{sec:fricas} --- its \verb|[]| is a theorem, not
a shrug.

The headline: \textbf{the free systems are not uniformly weaker here.}
Maxima solves the general case that SymPy cannot touch, and does so in
essentially the same closed form as Wolfram.  The relevant axis is not
commercial-vs-free but whether the system carries a confluent
hypergeometric ODE solver.  FriCAS sits outside that axis entirely: it
answers a different and more precise question, and the shape of its answer
turns out to encode the quantization of the energy (\S\ref{sec:fricas}).

\section{Wolfram Language}

\verb|DSolve| solves \eqref{eq:f1} outright.  Writing \verb|EE| for the
energy (plain \verb|E| is Euler's number and is \verb|Protected|):
\begin{lstlisting}
SetAttributes[{EE, m, beta1}, Constant];
ode = ( D[xi D[f1[xi], xi], xi]
        + ((1/2) EE xi - (1/4) m^2/xi + beta1) f1[xi] == 0 );
sol = DSolve[ode, f1, xi];
\end{lstlisting}
returning a single branch,
\begin{equation}\label{eq:wolfram}
f_1 = \xi^{m/2} e^{-i\sqrt{E/2}\,\xi}
      \left[ C_1\, U(a,\,m+1,\,i\sqrt{2E}\,\xi)
           + C_2\, L^{m}_{-a}(i\sqrt{2E}\,\xi) \right],
\qquad a = \tfrac{1}{2}\!\left(1+m+\frac{i\sqrt2\,\beta_1}{\sqrt E}\right).
\end{equation}
Setting $\beta_1=0$ first gives the Bessel reduction directly,
\begin{equation}
f_1 = C_1 J_{m/2}\!\left(\sqrt{E/2}\,\xi\right)
    + C_2 Y_{m/2}\!\left(\sqrt{E/2}\,\xi\right),
\end{equation}
which \verb|FullSimplify| verifies symbolically (returns \verb|True|).

Three practical traps, all of which cost real time:

\begin{enumerate}
\item \textbf{\texttt{.wls} line continuation.}  In a script file every
  top-level line that parses as a complete expression \emph{ends the
  statement}.  Writing \verb|ode = D[...]| on one line and
  \verb|+ [...] == 0| on the next silently assigns only the first line ---
  the entire bracketed potential term and the \verb|== 0| are discarded as
  separate throwaway expressions, and \verb|DSolve| then reports
  ``Equation or list of equations expected''.  Multi-line expressions must
  be parenthesised.

\item \textbf{\texttt{DSolve[ode, f1[xi], xi]} cannot be verified.}  That
  form returns a rule for the \emph{symbol} \verb|f1[xi]|, so \verb|f1'[xi]|
  and \verb|f1''[xi]| survive the replacement and the residual never
  closes.  Solving for the \emph{function}, \verb|DSolve[ode, f1, xi]|,
  returns \verb|f1 -> Function[{xi}, ...]| which substitutes into the
  derivatives as well.

\item \textbf{\texttt{FullSimplify} on \eqref{eq:wolfram} segfaults} the
  14.0 kernel when $m$ is symbolic (and then misreports the crash as a
  licence error).  The general solution was therefore verified numerically
  instead: the residual of \eqref{eq:f1} evaluates to $0$ at 40 digits for
  all of $E\in\{-1,3\}$, $\beta_1\in\{0,7/5\}$, $m\in\{0,2\}$.  The
  $\beta_1=0$ branch is small enough to verify symbolically.
\end{enumerate}

Also note that \verb|Print[InputForm[expr]]| prints the \emph{wrapper}, not
the formatted expression; \verb|ToString[expr, InputForm]| is what formats.
And \verb|TeXForm| of a \verb|Function| emits \verb|\unicode{f4a1}| for the
$\mapsto$ arrow, which will not compile --- extract \verb|f1[xi] /. sol[[i]]|
before asking for \TeX.

\section{Sage / Maxima}\label{sec:sage}

Plain \verb|desolve| fails: Maxima's \verb|ode2| does not recognise
\eqref{eq:f1} and raises \verb|NotImplementedError|.  With
\verb|contrib_ode=True| --- Maxima's \verb|contrib_ode| package, which
carries a Kovacic-style solver and hypergeometric pattern matching ---
it succeeds:
\begin{lstlisting}
sol = desolve(ode, f, ivar=xi, contrib_ode=True)
\end{lstlisting}
\begin{equation}\label{eq:maxima}
\begin{aligned}
f_1 &= \xi^{-m/2}\, e^{-z/2}
       \left[\, K_1\, M(a',\,1-m,\,z) \;+\; K_2\, U(a',\,1-m,\,z) \,\right], \\[4pt]
z &= -\sqrt{2}\,\sqrt{-E}\,\xi ,
\qquad
a' = -\,\frac{\sqrt{2}\,\sqrt{-E}\,\beta_1 + E\,(m-1)}{2E} .
\end{aligned}
\end{equation}
Passing \verb|ivar=xi| is mandatory; with four symbols in play Sage will
not guess the independent variable and raises ``Unable to determine
independent variable''.

\eqref{eq:maxima} is a correct solution, but it is the \emph{other}
Frobenius branch.  The indicial equation of \eqref{eq:f1x} at the regular
singular point $\xi=0$ has roots $\pm m/2$; Wolfram returns $\xi^{+m/2}$,
Maxima returns $\xi^{-m/2}$.  This matters twice over:

\begin{itemize}
\item The \verb|hypergeometric_M| half of \eqref{eq:maxima} is
  \textbf{undefined for positive integer $m$}, because its second parameter
  $1-m$ is then a non-positive integer and $M(a,b,z)$ has poles in $b$
  there.  Numerically this shows up as a division by zero at
  $m=2$ and (for $\beta_1\neq0$) $m=3$, while every non-integer $m$ and
  every case of the \verb|hypergeometric_U| half checks out to $10^{-31}$.
  Since $m$ is a magnetic quantum number, the degenerate case is precisely
  the physical one.
\item $\xi^{-m/2}$ is the solution that \emph{blows up} at the origin.  For
  \eqref{eq:f1} the admissible solution is the regular one, $\xi^{+m/2}$,
  so Wolfram's branch is the one to quote in the paper.
\end{itemize}

A smaller annoyance: differentiating Sage's \verb|hypergeometric_M|
symbolically raises
\begin{lstlisting}
ZeroDivisionError: symbolic division by zero
\end{lstlisting}
\noindent because its \verb|_derivative_| method forms $a/b$ with $b=1-m$.
The verification in \verb|f1-ode-sage.sage| is therefore done in mpmath.

For $\beta_1=0$ Maxima gives
$K_1 J_{-m/2}(-\sqrt{E/2}\,\xi) + K_2 Y_{-m/2}(-\sqrt{E/2}\,\xi)$ ---
again the $-m/2$ order, and with a negated argument, but a valid basis.

\section{SymPy}

SymPy solves \eqref{eq:f1} only when $\beta_1=0$.  \verb|classify_ode|
states the position exactly:

\begin{center}
\begin{tabular}{@{}ll@{}}
\toprule
case & hints offered \\
\midrule
$\beta_1=0$, any $m$ & \verb|factorable|, \verb|2nd_linear_bessel|, \verb|2nd_power_series_regular| \\
$\beta_1\neq0$       & \verb|factorable|, \verb|2nd_power_series_regular| \\
\bottomrule
\end{tabular}
\end{center}

\noindent With $\beta_1=0$ it returns the correct closed form, and for
symbolic $m$ at that:
\begin{lstlisting}
C1*besselj(m/2, sqrt(2)*sqrt(E)*xi/2) + C2*bessely(m/2, sqrt(2)*sqrt(E)*xi/2)
\end{lstlisting}
agreeing with Wolfram once $\sqrt2\sqrt E/2 = \sqrt{E/2}$ is recognised.

With $\beta_1\neq0$ the only applicable solver is Frobenius, so the answer
is a truncated series with an $O(\xi^6)$ tail rather than a closed form ---
and even that works only for numeric $m$; symbolic $m$ fails with
\verb|TypeError: Cannot convert symbols to int|.  (A related failure occurs
if $m$ is declared merely \verb|nonnegative|: the Frobenius code cannot
decide $-m/2 < m/2$ when $m$ might be $0$, and raises
\verb|TypeError: cannot determine truth value of Relational|.)

The cause is coverage, not cleverness.  SymPy 1.14 has \textbf{no Whittaker
functions at all} and no \verb|hyperu|/Kummer $U$; it carries \verb|hyper|
and \verb|meijerg| as generic objects but has no ODE solver that recognises
the confluent hypergeometric equation.  Wolfram reaches for
\verb|HypergeometricU| and \verb|LaguerreL| and Maxima for
\verb|hypergeometric_M|/\verb|_U|; SymPy has neither in usable form and
falls back to series.

SymPy also cannot \emph{verify} the closed form when handed one.  Given
\eqref{eq:wolfram} written as
\begin{lstlisting}
xi**(m/2)*exp(-I*k*xi)*hyper((a,),(m+1,),2*I*k*xi)
\end{lstlisting}
\noindent\verb|checkodesol| returns \verb|False| with a 711-operation residual: it
cannot collapse the contiguous-relation combination of \verb|hyper| terms.
The expression is nonetheless correct --- mpmath gives residuals of
$10^{-41}$ at 40 digits across $E\in\{0.7,3\}$, $\beta_1\in\{0,1.4\}$,
$m\in\{0,2,3\}$.  A \verb|False| from \verb|checkodesol| is a simplification
failure, not a refutation.


\section{FriCAS}\label{sec:fricas}

FriCAS returns
\begin{lstlisting}
[particular = 0, basis = []]
\end{lstlisting}
\noindent for \eqref{eq:f1} with symbolic parameters, and also for the
$\beta_1=0$ reduction.  Read carelessly this looks like the weakest result
of the four.  It is in fact the strongest statement any of them makes.

FriCAS's \verb|solve| for linear ODEs works in the \emph{Liouvillian}
class --- solutions built from rational functions by exponentiation,
integration and algebraic extension --- via the Kovacic algorithm, which
for second-order linear equations with rational coefficients is a
\emph{decision procedure}.  An empty basis therefore asserts that no
Liouvillian solution exists.  Bessel and confluent hypergeometric
functions are not Liouvillian, so there is nothing for it to return, and
saying so is the correct answer.

That the machinery is working, and that \verb|[]| is discriminating rather
than reflexive, is easy to confirm.  On Liouvillian problems it answers at
once:
\begin{center}
\begin{tabular}{@{}ll@{}}
\toprule
equation & basis returned \\
\midrule
$y''-y=0$                 & $e^{x},\; e^{-x}$ \\
$x^2y''+xy'-y=0$          & $(x^2-1)/x,\; (x^2+1)/x$ \\
$y''+x^{-1}y'=0$          & $1,\; \log x$ \\
\bottomrule
\end{tabular}
\end{center}
\noindent and FriCAS does carry Bessel functions (\verb|besselJ(1,2.0)|
evaluates to $0.5767248\ldots$) --- it simply will not manufacture them
from an ODE.

\subsection*{The $\beta_1=0$ case resolves by parity of $m$}

Bessel's equation of order $\nu$ has Liouvillian solutions if and only if
$\nu$ is a half-integer.  Here $\nu=m/2$, so the criterion is that $m$ be
\emph{odd}.  Setting $\beta_1=0$, $E=2$ and stepping $m$ reproduces this
exactly:

\begin{center}
\begin{tabular}{@{}llll@{}}
\toprule
$m$ & order $m/2$ & half-integer? & FriCAS \\
\midrule
1 & $1/2$ & yes & $e^{i\xi}/\sqrt\xi,\ e^{-i\xi}/\sqrt\xi$ \\
2 & $1$   & no  & \verb|[]| \\
3 & $3/2$ & yes & elementary (bulky) \\
4 & $2$   & no  & \verb|[]| \\
\bottomrule
\end{tabular}
\end{center}

\subsection*{With $\beta_1\neq0$ the criterion becomes quantization}

The confluent hypergeometric series terminates --- and the solution
becomes Liouvillian --- exactly when its first parameter is a non-positive
integer, i.e.\ when the Kummer function collapses to a Laguerre
polynomial.  For \eqref{eq:f1} with $E<0$ and $k=\sqrt{-E/2}$ that
condition reads
\begin{equation}\label{eq:quant}
\beta_1 = 2k\left(\frac{1+m}{2}+n_1\right), \qquad n_1=0,1,2,\dots
\end{equation}
which is precisely the bound-state quantization condition, $n_1$ the
parabolic quantum number.  Taking $E=-2$ so that $k=1$:

\begin{center}
\begin{tabular}{@{}lllp{0.42\textwidth}@{}}
\toprule
$m$ & $\beta_1$ & \eqref{eq:quant}? & FriCAS basis \\
\midrule
0 & 1 & yes, $n_1=0$ & $e^{-\xi}$,\; $\mathrm{Ei}(2\xi)e^{-\xi}$ \\
0 & 3 & yes, $n_1=1$ & $(2\xi-1)e^{-\xi}$,\; $\dots$ \\
0 & 2 & \textbf{no}  & \verb|[]| \\
2 & 3 & yes, $n_1=0$ & $\xi e^{-\xi}$,\; $\dots$ \\
2 & 2 & \textbf{no}  & \verb|[]| \\
\bottomrule
\end{tabular}
\end{center}

\noindent The first basis element is in each case $\xi^{m/2}e^{-k\xi}$
times the Laguerre polynomial of degree $n_1$; the second is the
reduction-of-order partner, Liouvillian but not elementary (hence the
exponential integral).  \textbf{FriCAS's Liouvillian decision procedure is
detecting the energy quantization.}  Where Wolfram and Maxima hand back a
closed form valid for all $\beta_1$, FriCAS partitions the $\beta_1$ axis
into the physically admissible points and the rest.  Neither answer
subsumes the other, and for this problem FriCAS's is arguably the more
informative.

\subsection*{Installation note}

FriCAS is not packaged for this machine and needs no root:
\begin{lstlisting}
mamba create -n fricas -c conda-forge fricas
mamba install -n fricas -c conda-forge 'sbcl=2.4.8'
\end{lstlisting}
\noindent The second command is required.  The conda \verb|fricas| 1.3.13
core is built against SBCL 2.4.8, the solver installs a newer SBCL, and the
mismatch aborts at startup with
\begin{lstlisting}
GC strategy mismatch: runtime uses 1, core built for 38
\end{lstlisting}
\noindent Also, \verb|-nosman| must be the \emph{first}
argument; otherwise FriCAS starts its session manager, which wants a tty
and hangs on piped stdin.

\section{Files}

\begin{tabular}{@{}ll@{}}
\toprule
\verb|cas/f1-ode.wls|       & Wolfram script; \verb|DSolve|, verification, Whittaker cross-check \\
\verb|cas/f1-ode-sympy.py|  & SymPy script; \verb|classify_ode| evidence, mpmath verification \\
\verb|cas/f1-ode-sage.sage| & Sage/Maxima script; \verb|contrib_ode|, branch degeneracy check \\
\verb|cas/f1-ode-fricas.input| & FriCAS script; Liouvillian decision, quantization scan \\
\verb|cas/cas-comparison.tex| & this note \\
\bottomrule
\end{tabular}

\bigskip
\noindent Run them with, respectively:
\begin{lstlisting}
/opt/Wolfram/WolframScript/bin/wolframscript -file f1-ode.wls
python f1-ode-sympy.py
sage f1-ode-sage.sage
fricas -nosman < f1-ode-fricas.input
\end{lstlisting}

\vspace{2em}
\noindent\rule{0.3\textwidth}{0.4pt}\\
{\footnotesize /claude-opus-5}

\end{document}
